\documentclass[10pt,a4paper]{amsart}
\usepackage{amsmath,amssymb,amsthm,graphicx}
\usepackage{verbatim}
\newenvironment{reference}{\comment}{\endcomment}
%\newenvironment{reference}{}{}
\newenvironment{slogan}{\comment}{\endcomment}
\newenvironment{history}{\comment}{\endcomment}

% For commutative diagrams we use Xy-pic
\usepackage[all]{xy}

% We use 2cell for 2-commutative diagrams.
\xyoption{2cell}
% Theorem environments.
%
\theoremstyle{plain}
\newtheorem{theorem}[subsection]{Theorem}
\newtheorem{proposition}[subsection]{Proposition}
\newtheorem{lemma}[subsection]{Lemma}

\theoremstyle{definition}
\newtheorem{definition}[subsection]{Definition}
\newtheorem{example}[subsection]{Example}
\newtheorem{exercise}[subsection]{Exercise}
\newtheorem{situation}[subsection]{Situation}

\theoremstyle{remark}
\newtheorem{remark}[subsection]{Remark}
\newtheorem{remarks}[subsection]{Remarks}

\numberwithin{equation}{subsection}

% Macros
%
\def\lim{\mathop{\mathrm{lim}}\nolimits}
\def\colim{\mathop{\mathrm{colim}}\nolimits}
\def\Spec{\mathop{\mathrm{Spec}}}
\def\Hom{\mathop{\mathrm{Hom}}\nolimits}
\def\Ext{\mathop{\mathrm{Ext}}\nolimits}
\def\SheafHom{\mathop{\mathcal{H}\!\mathit{om}}\nolimits}
\def\SheafExt{\mathop{\mathcal{E}\!\mathit{xt}}\nolimits}
\def\Sch{\mathit{Sch}}
\def\Mor{\mathop{\mathrm{Mor}}\nolimits}
\def\Ob{\mathop{\mathrm{Ob}}\nolimits}
\def\Sh{\mathop{\mathit{Sh}}\nolimits}
\def\NL{\mathop{N\!L}\nolimits}
\def\CH{\mathop{\mathrm{CH}}\nolimits}
\def\proetale{{pro\text{-}\acute{e}tale}}
\def\etale{{\acute{e}tale}}
\def\QCoh{\mathit{QCoh}}
\def\Ker{\mathop{\mathrm{Ker}}}
\def\Im{\mathop{\mathrm{Im}}}
\def\Coker{\mathop{\mathrm{Coker}}}
\def\Coim{\mathop{\mathrm{Coim}}}

% Boxtimes
%
\DeclareMathSymbol{\boxtimes}{\mathbin}{AMSa}{"02}

%
% Macros for moduli stacks/spaces
%
\def\QCohstack{\mathcal{QC}\!\mathit{oh}}
\def\Cohstack{\mathcal{C}\!\mathit{oh}}
\def\Spacesstack{\mathcal{S}\!\mathit{paces}}
\def\Quotfunctor{\mathrm{Quot}}
\def\Hilbfunctor{\mathrm{Hilb}}
\def\Curvesstack{\mathcal{C}\!\mathit{urves}}
\def\Polarizedstack{\mathcal{P}\!\mathit{olarized}}
\def\Complexesstack{\mathcal{C}\!\mathit{omplexes}}
% \Pic is the operator that assigns to X its picard group, usage \Pic(X)
% \Picardstack_{X/B} denotes the Picard stack of X over B
% \Picardfunctor_{X/B} denotes the Picard functor of X over B
\def\Pic{\mathop{\mathrm{Pic}}\nolimits}
\def\Picardstack{\mathcal{P}\!\mathit{ic}}
\def\Picardfunctor{\mathrm{Pic}}
\def\Deformationcategory{\mathcal{D}\!\mathit{ef}}
\usepackage{cleveref}
\crefname{originalreference}{original reference}{original references}
\Crefname{originalreference}{Original reference}{Original references}
\makeatletter
\newlabel{lemma-get-triangle}{{09P6}{source}{}{sourcepointer-1}{}}
\newlabel{lemma-get-triangle@cref}{{[originalreference][0][]09P6}{[1][0][]source}{}{}{}}
\makeatother
\crefname{originalreference}{original reference}{original references}
\Crefname{originalreference}{Original reference}{Original references}
\makeatletter
\newlabel{situation-ABC}{{09QJ}{source}{}{sourcepointer-2}{}}
\newlabel{situation-ABC@cref}{{[originalreference][0][]09QJ}{[1][0][]source}{}{}{}}
\makeatother
\title{The rotated admissible triangle is isomorphic to the cone triangle of the rotated differential}
\author{The Stacks Project Authors}
\date{stacks-project:dga.tex:local-snapshot-20261012; GNU FDL 1.2 or later}
\begin{document}
\maketitle
\noindent\textit{Editorial note.} All delivered mathematical text is the author's own, copied from the pinned Stacks Project chapter file sources/stacks-dga/pinned-original-source.tex; no statement or proof step was rewritten, reordered or removed.\par\medskip
\noindent\textit{Editorial note.} The counted proof is the author's complete argument: the splitting data and the two comparison maps, the differential closedness computations with the d(tilde-alpha) and d(tilde-p) identities, the diagram fit, and both composition identities collapsing to the respective identity morphisms are all retained verbatim (the author marks the diagram fit as immediate; the closedness and composition computations that carry the proof are displayed in full).\par\medskip
\noindent\textit{Editorial note.} The label lemma-cone-rotate-isom names the counted statement, and the proof environment immediately following it in the source is delivered verbatim as its proof.\par\medskip
\noindent\textit{Editorial note.} The counted statement is carried by the author's own lemma environment in the chapter source.\par\medskip
\noindent\textit{Allowed standard prerequisites (editorial).}\begin{itemize}
\item Identities for the splitting data of an admissible sequence. Same chapter, Lemma lemma-get-triangle. Supplies the relations delta[-1] = tilde-alpha d(tilde-beta) = -d(tilde-alpha) tilde-beta and delta[-1] = tilde-i d(tilde-p) together with equation (6), which drive every differential check and the composition computation.
\item Cones of morphisms in the differential graded category. Same chapter, Definition definition-cone and Situation situation-ABC. Fixes the cone c(delta[-1]) with its structure maps i and p whose triangle is compared with the rotated admissible triangle.
\end{itemize}\medskip
\ifdefined\hypertarget\hypertarget{sourcepointer-1}{}\fi
\noindent\textit{Source reference.} Original reference 09P6: sources/stacks-dga/pinned-original-source.tex, source lines 3708--3708. Prerequisite pointer invoked at proof line(s) 4207; it names a standard fact of the same chapter that the counted argument uses, and every step of the counted proof built on it is retained verbatim in the delivered body.\par\medskip
\ifdefined\hypertarget\hypertarget{sourcepointer-2}{}\fi
\noindent\textit{Source reference.} Original reference 09QJ: sources/stacks-dga/pinned-original-source.tex, source lines 3762--3762. Prerequisite pointer invoked at proof line(s) 4168; it names a standard fact of the same chapter that the counted argument uses, and every step of the counted proof built on it is retained verbatim in the delivered body.\par\medskip
\begin{lemma}
\label{lemma-cone-rotate-isom}
In Situation \ref{situation-ABC} let $\alpha : x \to y$ and $\beta : y \to z$
define an admissible short exact sequence
$$
\xymatrix{
x \ar[r] &
y\ar[r] &
z
}
$$
in $\text{Comp}(\mathcal{A})$. Let $(x, y, z, \alpha, \beta, \delta)$
be the associated triangle in $K(\mathcal{A})$. Then, the triangles
$$
(z[-1], x, y, \delta[-1], \alpha, \beta)
\quad\text{and}\quad
(z[-1], x, c(\delta[-1]), \delta[-1], i, p)
$$
are isomorphic. 
\end{lemma}

\begin{proof}
We have a diagram of the form
$$
\xymatrix{
z[-1]\ar[r]^{\delta[-1]}\ar[d]^1 &
x\ar@<0.5ex>[r]^{\alpha}\ar[d]^1 &
y\ar@<0.5ex>[r]^{\beta}\ar@{.>}[d]\ar@<0.5ex>[l]^{\tilde{\alpha}} &
z\ar[d]^1\ar@<0.5ex>[l]^{\tilde\beta} \\
z[-1] \ar[r]^{\delta[-1]} &
x\ar@<0.5ex>[r]^i &
c(\delta[-1]) \ar@<0.5ex>[r]^p\ar@<0.5ex>[l]^{\tilde i} &
z\ar@<0.5ex>[l]^{\tilde p}
}
$$
with splittings to $\alpha, \beta, i$, and $p$ given by
$\tilde{\alpha}, \tilde{\beta}, \tilde{i},$ and $\tilde{p}$ respectively.
Define a morphism $y \to c(\delta[-1])$ by
$i\tilde{\alpha} + \tilde{p}\beta$ and a morphism
$c(\delta[-1]) \to y$ by $\alpha \tilde{i} + \tilde{\beta} p$.
Let us first check that these define morphisms in $\text{Comp}(\mathcal{A})$.
We remark that by identities from Lemma \ref{lemma-get-triangle},
we have the relation
$\delta[-1] = \tilde{\alpha}d(\tilde{\beta}) = -d(\tilde{\alpha})\tilde{\beta}$
and the relation $\delta[-1] = \tilde{i}d(\tilde{p})$. Then
\begin{align*}
d(\tilde{\alpha})
& =
d(\tilde{\alpha})\tilde{\beta}\beta \\
& =
-\delta[-1]\beta
\end{align*}
where we have used equation (6) of
Lemma \ref{lemma-get-triangle} for the first equality and
the preceding remark for the second. Similarly, we obtain
$d(\tilde{p}) = i\delta[-1]$. Hence
\begin{align*}
d(i\tilde{\alpha} + \tilde{p}\beta)
& =
d(i)\tilde{\alpha} + id(\tilde{\alpha}) +
d(\tilde{p})\beta + \tilde{p}d(\beta) \\
& =
id(\tilde{\alpha}) + d(\tilde{p})\beta \\
& =
-i\delta[-1]\beta + i\delta[-1]\beta \\
& =
0
\end{align*}
so $i\tilde{\alpha} + \tilde{p}\beta$ is indeed a morphism of
$\text{Comp}(\mathcal{A})$. By a similar calculation,
$\alpha \tilde{i} + \tilde{\beta} p$ is also a morphism of
$\text{Comp}(\mathcal{A})$. It is immediate that these morphisms
fit in the commutative diagram. We compute:
\begin{align*}
(i\tilde{\alpha} + \tilde{p}\beta)(\alpha \tilde{i} + \tilde{\beta} p)
& =
i\tilde{\alpha}\alpha\tilde{i} + i\tilde{\alpha}\tilde{\beta}p
+ \tilde{p}\beta\alpha\tilde{i} + \tilde{p}\beta\tilde{\beta}p \\
& =
i\tilde{i} + \tilde{p}p \\
& =
1_{c(\delta[-1])}
\end{align*}
where we have freely used the identities of
Lemma \ref{lemma-get-triangle}. Similarly, we compute
$(\alpha \tilde{i} + \tilde{\beta} p)(i\tilde{\alpha} + \tilde{p}\beta) = 1_y$,
so we conclude $y \cong c(\delta[-1])$. Hence, the two triangles in question
are isomorphic.
\end{proof}

\end{document}
